% =============================================================================
% Appendix B — Glossary
% =============================================================================
\chapter{Glossary}\label{app:glossary}

This glossary defines the key terms, acronyms, and technologies referenced
throughout this book. Entries are grouped alphabetically. Cross-references to
related terms appear in \textit{italics}.

%% ---------------------------------------------------------------------------
\section*{A}
%% ---------------------------------------------------------------------------

\begin{description}

\item[\textbf{ACL (Access Control List)}]
A fine-grained permission mechanism that extends traditional UNIX
owner/group/other modes. POSIX ACLs attach per-user or per-group entries to
individual files or directories. NFSv4 ACLs provide richer semantics
including inheritance and deny rules.

\item[\textbf{AFF (All Flash FAS)}]
NetApp hardware platform built exclusively with solid-state drives. AFF
systems run \textit{ONTAP} and are optimised for low-latency, high-IOPS
workloads such as databases and virtualised environments.

\item[\textbf{ALUA (Asymmetric Logical Unit Access)}]
A SCSI standard that allows a storage array to advertise preferred and
non-preferred paths to a \textit{LUN}. The host multipath driver uses this
information to route I/O through the optimal controller.

\item[\textbf{ANA (Asymmetric Namespace Access)}]
The NVMe equivalent of \textit{ALUA}. ANA enables multipath-aware hosts to
identify the optimal controller for a given NVMe namespace.

\item[\textbf{ARC (Adaptive Replacement Cache)}]
The primary read cache in \textit{ZFS}, held in RAM. ARC uses a
frequency-and-recency algorithm to keep hot blocks cached. It dynamically
grows and shrinks in response to system memory pressure.

\item[\textbf{ARP (Address Resolution Protocol)}]
A Layer-2 protocol that resolves IPv4 addresses to MAC addresses on a local
network segment. Relevant to storage networking when iSCSI or NFS traffic
traverses Ethernet.

\end{description}

%% ---------------------------------------------------------------------------
\section*{B}
%% ---------------------------------------------------------------------------

\begin{description}

\item[\textbf{BBU (Battery Backup Unit)}]
A battery module attached to a RAID controller or NVRAM card that protects
in-flight write data during a power failure, enabling the controller to
safely use a write-back cache policy.

\item[\textbf{BFQ (Budget Fair Queuing)}]
A Linux I/O scheduler designed for interactive and latency-sensitive
workloads. BFQ allocates per-process I/O budgets to provide fairness and
low latency, particularly on rotational media.

\item[\textbf{BlueXP}]
NetApp's unified control plane for managing hybrid and multi-cloud storage
environments. BlueXP provides discovery, provisioning, data protection, and
governance across on-premises ONTAP, \textit{CVO}, \textit{FSx for ONTAP},
\textit{GCNV}, and \textit{StorageGRID} deployments.

\item[\textbf{Btrfs (B-tree Filesystem)}]
A Linux copy-on-write filesystem supporting integrated volume management,
snapshots, checksumming, compression, and multi-device spanning. Btrfs aims
to be a general-purpose next-generation filesystem for Linux.

\end{description}

%% ---------------------------------------------------------------------------
\section*{C}
%% ---------------------------------------------------------------------------

\begin{description}

\item[\textbf{CAP Theorem}]
A distributed systems principle stating that a networked data store can
simultaneously guarantee at most two of: Consistency, Availability, and
Partition tolerance. Storage architects use it to reason about trade-offs
in distributed and replicated systems.

\item[\textbf{CBT (Changed Block Tracking)}]
A mechanism that records which disk blocks have changed since a previous
point in time. CBT accelerates incremental backup and replication by
allowing the backup engine to read only modified blocks.

\item[\textbf{CDM (Copy Data Management)}]
A strategy and toolset for reducing the proliferation of redundant data
copies (development, test, analytics) by using space-efficient clones and
snapshots instead of full physical copies.

\item[\textbf{CephFS}]
The POSIX-compliant distributed filesystem built on top of the \textit{RADOS}
object store in Ceph. CephFS uses one or more Metadata Servers
(\textit{MDS}) to manage the namespace and delegates data storage to
\textit{OSDs}.

\item[\textbf{CHAP (Challenge-Handshake Authentication Protocol)}]
An authentication scheme used in \textit{iSCSI} to verify the identity of
initiators and targets. Mutual CHAP authenticates both sides of the
connection.

\item[\textbf{CIFS (Common Internet File System)}]
The original name for the SMB file-sharing protocol as extended by
Microsoft. Modern implementations use \textit{SMB} 2.x/3.x, but the term
CIFS persists in Linux mount helpers (\texttt{mount.cifs}).

\item[\textbf{Cloud Volumes ONTAP (CVO)}]
A software-defined storage instance that runs the NetApp \textit{ONTAP}
operating system on public cloud infrastructure (AWS, Azure, GCP). CVO
provides enterprise data services---snapshots, replication, tiering---in the
cloud.

\item[\textbf{CoW (Copy-on-Write)}]
A technique where data blocks are never overwritten in place. Instead, a
modified copy is written to a new location and the metadata pointer is
atomically updated. Used by \textit{ZFS}, \textit{Btrfs}, and
\textit{WAFL}.

\item[\textbf{CRUSH (Controlled Replication Under Scalable Hashing)}]
The deterministic data placement algorithm in \textit{Ceph}. CRUSH maps
objects to \textit{OSDs} based on a hierarchical cluster map without
requiring a central lookup table.

\item[\textbf{CSI (Container Storage Interface)}]
A standard API that allows Kubernetes and other container orchestrators to
provision, attach, mount, snapshot, and resize persistent storage volumes
from any compliant driver. See also \textit{Trident}.

\item[\textbf{CVO}]
See \textit{Cloud Volumes ONTAP}.

\end{description}

%% ---------------------------------------------------------------------------
\section*{D}
%% ---------------------------------------------------------------------------

\begin{description}

\item[\textbf{DAID (Data-At-Rest Encryption Identifier)}]
An identifier associated with a self-encrypting drive's encryption domain,
used when managing keys for \textit{SED} and \textit{NSE} devices within a
storage system.

\item[\textbf{DM-Multipath}]
The Linux device-mapper multipath framework. DM-Multipath aggregates
multiple physical paths between a host and a storage device into a single
logical device, providing redundancy and optional load balancing.

\item[\textbf{DORA (Discover, Offer, Request, Acknowledge)}]
The four-step handshake in DHCP used to assign IP addresses. In storage
contexts, DORA is relevant to iSCSI boot and PXE environments where
initiators obtain their network identity via DHCP.

\item[\textbf{dRAID (Distributed RAID)}]
A \textit{ZFS} pool topology that distributes parity and spare capacity
across all drives in a group, dramatically accelerating rebuild times
compared to traditional RAID-Z by parallelising reconstruction I/O.

\item[\textbf{DRBD (Distributed Replicated Block Device)}]
A Linux kernel module that mirrors block devices in real time over a
network. DRBD is often called ``network RAID-1'' and is used for
high-availability clustering.

\item[\textbf{DWPD (Drive Writes Per Day)}]
An SSD endurance metric expressing how many times the entire drive capacity
can be overwritten per day over the warranty period before wear-out. See
also \textit{TBW}.

\end{description}

%% ---------------------------------------------------------------------------
\section*{E}
%% ---------------------------------------------------------------------------

\begin{description}

\item[\textbf{Erasure Coding}]
A data protection scheme that splits an object into $k$ data fragments and
computes $m$ parity (coding) fragments such that any $k$ of the $k+m$
fragments suffice to reconstruct the original data. More space-efficient
than replication for cold and warm data.

\item[\textbf{ext4 (Fourth Extended Filesystem)}]
The mature, stable default filesystem on many Linux distributions.\ ext4
supports volumes up to 1\,EiB, extents-based allocation, delayed
allocation, and journaling.

\end{description}

%% ---------------------------------------------------------------------------
\section*{F}
%% ---------------------------------------------------------------------------

\begin{description}

\item[\textbf{FabricPool}]
A NetApp \textit{ONTAP} feature that automatically tiers cold data blocks
from a high-performance local tier (SSD aggregate) to a low-cost object
store (S3, Azure Blob, Google Cloud Storage, or \textit{StorageGRID}).

\item[\textbf{FAS (Fabric-Attached Storage)}]
NetApp's hybrid (HDD\,+\,SSD) storage platform running \textit{ONTAP}.
FAS systems support both \textit{SAN} and \textit{NAS} protocols.

\item[\textbf{FC (Fibre Channel)}]
A high-speed serial transport protocol designed for storage networking.\ FC
provides lossless, in-order delivery and is the traditional backbone of
enterprise \textit{SAN} environments. Common speeds include 32\,Gbps and
64\,Gbps.

\item[\textbf{FlexCache}]
An ONTAP feature that creates a sparse, read-write caching volume at a
remote site, reducing latency for frequently accessed data from a distant
origin \textit{FlexVol} or \textit{FlexGroup}.

\item[\textbf{FlexClone}]
A NetApp technology that creates writable, space-efficient clones of
\textit{FlexVol} volumes or LUNs using shared block pointers. Only changed
blocks consume additional space.

\item[\textbf{FlexGroup}]
An ONTAP scale-out NAS container that distributes a single namespace across
multiple constituent \textit{FlexVol} volumes and nodes, delivering high
throughput and large capacity for unstructured data workloads.

\item[\textbf{FlexVol}]
The fundamental logical volume in \textit{ONTAP}. A FlexVol resides within
an aggregate and can host NFS exports, CIFS shares, or \textit{LUNs}.

\item[\textbf{FLOGI (Fabric Login)}]
The initial Fibre Channel login process in which a node registers with the
fabric switch and obtains an N\_Port ID.\ FLOGI is analogous to DHCP in IP
networking.

\item[\textbf{FPolicy}]
An ONTAP file-access notification framework that intercepts file operations
(open, create, rename, delete) and forwards events to an external server
for real-time screening, auditing, or data classification.

\item[\textbf{FSx for ONTAP}]
A fully managed AWS service that provides NetApp \textit{ONTAP} file systems
in the cloud. FSx for ONTAP supports NFS, SMB, and iSCSI with features
such as snapshots, \textit{FlexClone}, and \textit{SnapMirror}.

\end{description}

%% ---------------------------------------------------------------------------
\section*{G}
%% ---------------------------------------------------------------------------

\begin{description}

\item[\textbf{GCNV (Google Cloud NetApp Volumes)}]
A Google Cloud--native service built on NetApp \textit{ONTAP} technology that
delivers fully managed NFS and SMB file storage with enterprise data
services.

\item[\textbf{GDPR (General Data Protection Regulation)}]
European Union regulation governing the processing and storage of personal
data. GDPR mandates data minimisation, retention limits, and the right to
erasure---all of which have direct implications for storage architecture
and lifecycle management.

\item[\textbf{GDS (GPUDirect Storage)}]
See \textit{GPUDirect Storage}.

\item[\textbf{GFS2 (Global File System 2)}]
A shared-disk cluster filesystem for Linux that uses distributed locking
(\texttt{dlm}) to allow multiple nodes to mount and write to the same
block device concurrently.

\item[\textbf{GPFS (General Parallel File System)}]
IBM's high-performance parallel filesystem (now marketed as IBM Storage
Scale). GPFS stripes data across many disks and servers and is widely
deployed in HPC and AI environments.

\item[\textbf{GPUDirect Storage}]
An NVIDIA technology that enables a direct data path between storage and GPU
memory, bypassing the CPU and system memory entirely. GPUDirect Storage
reduces latency and CPU overhead for AI/ML training and inference
pipelines.

\end{description}

%% ---------------------------------------------------------------------------
\section*{H}
%% ---------------------------------------------------------------------------

\begin{description}

\item[\textbf{HBA (Host Bus Adapter)}]
A hardware interface card that connects a host to a storage network.\ FC
HBAs connect servers to Fibre Channel SANs; iSCSI HBAs (or TOE cards)
offload TCP/IP processing for IP-based SANs.

\item[\textbf{HCI (Hyper-Converged Infrastructure)}]
An architecture that combines compute, storage, and networking into a single
appliance or software stack, managed as a unified platform. Examples
include VMware vSAN, Nutanix, and Azure Stack HCI.

\item[\textbf{HSM (Hierarchical Storage Management)}]
A policy-driven system that automatically migrates data between storage
tiers (e.g., SSD to HDD to tape or cloud) based on access frequency,
age, or other criteria.

\end{description}

%% ---------------------------------------------------------------------------
\section*{I}
%% ---------------------------------------------------------------------------

\begin{description}

\item[\textbf{IOPS (Input/Output Operations Per Second)}]
The primary throughput metric for storage devices handling small-block
random workloads. IOPS must always be qualified by block size, read/write
mix, and queue depth to be meaningful.

\item[\textbf{iSCSI (Internet Small Computer Systems Interface)}]
A transport protocol that encapsulates \textit{SCSI} commands inside TCP/IP
packets, enabling block-level storage access over standard Ethernet
networks.

\item[\textbf{iSNS (Internet Storage Name Service)}]
A discovery protocol for iSCSI and Fibre Channel over Ethernet environments.
iSNS provides automated discovery, registration, and management of storage
nodes on an IP network.

\end{description}

%% ---------------------------------------------------------------------------
\section*{K}
%% ---------------------------------------------------------------------------

\begin{description}

\item[\textbf{Kerberos}]
A network authentication protocol based on symmetric-key cryptography and a
trusted third-party ticket-granting service. In storage, Kerberos secures
NFSv4 (\texttt{sec=krb5/krb5i/krb5p}) and SMB sessions.

\item[\textbf{KMIP (Key Management Interoperability Protocol)}]
An OASIS standard for communication between key management servers and
encryption clients. Storage systems use KMIP to manage encryption keys for
\textit{SED}, \textit{NVE}, \textit{NAE}, and \textit{LUKS} volumes.

\end{description}

%% ---------------------------------------------------------------------------
\section*{L}
%% ---------------------------------------------------------------------------

\begin{description}

\item[\textbf{L2ARC (Level-2 Adaptive Replacement Cache)}]
A secondary read cache in \textit{ZFS} stored on a fast device (typically
NVMe or SSD). The L2ARC extends the in-memory \textit{ARC} for
working sets that exceed available RAM.

\item[\textbf{LIF (Logical Interface)}]
A virtual network endpoint in \textit{ONTAP} that can migrate between
physical ports and nodes. LIFs provide protocol-specific access (NFS,
CIFS, iSCSI, FC) and enable non-disruptive operations.

\item[\textbf{Little's Law}]
A queueing-theory result relating average queue length ($L$), average
arrival rate ($\lambda$), and average wait time ($W$): $L = \lambda W$.
Widely used to model storage I/O latency and queue depth.

\item[\textbf{LUKS (Linux Unified Key Setup)}]
The standard on-disk format for full-disk encryption on Linux. LUKS stores
key material and cipher metadata in a header, managed by
\texttt{cryptsetup}. LUKS2 adds Argon2 key derivation and authenticated
encryption.

\item[\textbf{LUN (Logical Unit Number)}]
A logical address that identifies a specific storage volume or partition
presented over a \textit{SAN} (FC, iSCSI, or NVMe-oF). From the host
perspective, a LUN appears as a local block device.

\item[\textbf{LVM (Logical Volume Manager)}]
A Linux subsystem that provides an abstraction layer between physical
disks/partitions and filesystems. LVM enables online resizing, snapshots,
thin provisioning, and striping.

\item[\textbf{Lustre}]
A POSIX-compliant parallel distributed filesystem designed for HPC clusters.
Lustre separates metadata (managed by \textit{MDS} nodes) from data
(stored on \textit{OST} targets) to deliver massive aggregate throughput.

\end{description}

%% ---------------------------------------------------------------------------
\section*{M}
%% ---------------------------------------------------------------------------

\begin{description}

\item[\textbf{MAV (Multi-Admin Verification)}]
An ONTAP security feature that requires approval from multiple designated
administrators before critical operations (e.g., volume deletion or
SnapLock modification) can be executed.

\item[\textbf{MD-RAID}]
The Linux kernel's software RAID implementation managed by \texttt{mdadm}.
MD-RAID supports levels 0, 1, 4, 5, 6, and 10, with online reshape and
hot-spare capabilities.

\item[\textbf{MDS (Metadata Server)}]
A server or daemon responsible for managing filesystem namespace operations
(directory listing, file creation, permission checks). Used in
\textit{CephFS}, \textit{Lustre}, and other parallel filesystems.

\item[\textbf{MetroCluster}]
A NetApp high-availability and disaster-recovery solution that synchronously
mirrors data between two geographically separated ONTAP clusters using
\textit{SyncMirror} over dedicated links.

\item[\textbf{MinIO}]
A high-performance, Kubernetes-native object storage server that implements
the S3 API. MinIO is commonly deployed for on-premises or hybrid-cloud
object storage.

\end{description}

%% ---------------------------------------------------------------------------
\section*{N}
%% ---------------------------------------------------------------------------

\begin{description}

\item[\textbf{NAE (NetApp Aggregate Encryption)}]
An ONTAP software encryption feature that encrypts all data within an
aggregate using AES-256, providing volume-level encryption without
requiring \textit{SED} drives. See also \textit{NVE}.

\item[\textbf{NAS (Network-Attached Storage)}]
A file-level storage architecture in which a dedicated appliance or server
shares filesystems over a network using protocols such as \textit{NFS} or
\textit{SMB/CIFS}. Clients mount remote exports as local directories.

\item[\textbf{NFS (Network File System)}]
A distributed filesystem protocol that allows clients to access files over a
network as though they were local. NFSv4.x adds stateful operation,
compound RPCs, \textit{pNFS} extensions, and strong security via
\textit{Kerberos}.

\item[\textbf{NIS2 (Network and Information Security Directive 2)}]
An EU directive that imposes cybersecurity risk-management and
incident-reporting obligations on operators of essential and important
services, including data centre and cloud providers.

\item[\textbf{NSE (NetApp Storage Encryption)}]
Hardware-based encryption at rest using \textit{SED} (self-encrypting
drives) managed by ONTAP. NSE protects against physical drive theft without
software overhead.

\item[\textbf{NVE (NetApp Volume Encryption)}]
A software-based ONTAP encryption feature that encrypts data at the volume
(FlexVol/FlexGroup) level using AES-256. NVE can layer on top of NSE for
double encryption.

\item[\textbf{NVMe (Non-Volatile Memory Express)}]
A host controller interface and storage protocol optimised for flash and
next-generation non-volatile memory. NVMe uses PCIe lanes and deep
parallel queues (up to 64\,K queues, each 64\,K entries deep) to minimise
latency.

\item[\textbf{NVMe-oF (NVMe over Fabrics)}]
An extension of the NVMe protocol that runs over network fabrics such as
\textit{RoCE}, Fibre Channel (\textit{FC-NVMe}), or TCP. NVMe-oF brings
NVMe-class latency to remote storage.

\item[\textbf{NVRAM (Non-Volatile Random-Access Memory)}]
Battery- or capacitor-backed memory used in storage controllers to
journal incoming writes before they are committed to persistent media.
NVRAM enables fast write acknowledgement while guaranteeing durability.

\end{description}

%% ---------------------------------------------------------------------------
\section*{O}
%% ---------------------------------------------------------------------------

\begin{description}

\item[\textbf{OCFS2 (Oracle Cluster File System 2)}]
A shared-disk cluster filesystem for Linux developed by Oracle. OCFS2 allows
multiple nodes to concurrently read and write to shared storage, often used
with Oracle RAC databases.

\item[\textbf{OKM (Onboard Key Manager)}]
An ONTAP component that stores encryption keys locally on the storage
cluster using a secure key hierarchy, as an alternative to an external
\textit{KMIP} server.

\item[\textbf{ONTAP}]
NetApp's storage operating system powering \textit{FAS}, \textit{AFF},
\textit{CVO}, \textit{FSx for ONTAP}, and \textit{GCNV}. ONTAP provides
unified (SAN + NAS) multiprotocol access, \textit{WAFL},
\textit{SnapMirror}, and comprehensive data services.

\item[\textbf{OSD (Object Storage Daemon)}]
A Ceph daemon responsible for storing data objects, handling replication,
recovery, rebalancing, and scrubbing on a single storage device. Each OSD
typically maps to one physical disk.

\item[\textbf{OST (Object Storage Target)}]
A \textit{Lustre} component that stores file data objects. Lustre stripes
file content across multiple OSTs to maximise aggregate throughput.

\end{description}

%% ---------------------------------------------------------------------------
\section*{P}
%% ---------------------------------------------------------------------------

\begin{description}

\item[\textbf{pNFS (Parallel NFS)}]
An NFSv4.1 extension that allows clients to perform direct, parallel I/O to
multiple data servers. The metadata server provides layout maps so clients
can bypass the single NFS server bottleneck.

\item[\textbf{PV (Persistent Volume)}]
A Kubernetes object representing a piece of storage in the cluster.\ PVs
are provisioned by administrators or dynamically via a \textit{CSI} driver
and bound to \textit{PVCs} by matching capacity and access modes.

\item[\textbf{PVC (Persistent Volume Claim)}]
A Kubernetes resource that requests storage of a specified size and access
mode. The cluster binds the PVC to a matching \textit{PV}.

\end{description}

%% ---------------------------------------------------------------------------
\section*{Q}
%% ---------------------------------------------------------------------------

\begin{description}

\item[\textbf{QoS (Quality of Service)}]
A set of mechanisms that guarantee minimum throughput and/or cap maximum IOPS
and bandwidth for a storage workload. QoS policies prevent noisy-neighbour
effects in multi-tenant environments.

\end{description}

%% ---------------------------------------------------------------------------
\section*{R}
%% ---------------------------------------------------------------------------

\begin{description}

\item[\textbf{RADOS (Reliable Autonomic Distributed Object Store)}]
The foundational object storage layer of \textit{Ceph}. RADOS distributes
objects across \textit{OSDs} using the \textit{CRUSH} algorithm and
provides self-healing replication or erasure coding.

\item[\textbf{RAID (Redundant Array of Independent Disks)}]
A family of techniques that combine multiple physical drives to improve
performance, capacity, or fault tolerance. Common levels include RAID-0
(striping), RAID-1 (mirroring), RAID-5 (single parity), and RAID-6
(double parity).

\item[\textbf{RAID-DP}]
NetApp's proprietary double-parity RAID implementation in \textit{ONTAP},
functionally equivalent to RAID-6 but optimised for the \textit{WAFL}
filesystem layout.

\item[\textbf{RAID-TEC}]
NetApp's triple-parity RAID for \textit{ONTAP}, capable of tolerating three
simultaneous disk failures within a RAID group. Recommended for
large-capacity HDDs.

\item[\textbf{RAID-Z}]
The \textit{ZFS} family of distributed-parity schemes. RAID-Z1, Z2, and Z3
tolerate one, two, or three device failures respectively. RAID-Z avoids
the write hole by using variable-width stripes.

\item[\textbf{RBD (RADOS Block Device)}]
A block storage interface layered on \textit{RADOS} in Ceph. RBD images
are thin-provisioned, snapshotable, and can be mapped as local block
devices or consumed by virtual machines and Kubernetes pods.

\item[\textbf{RDMA (Remote Direct Memory Access)}]
A networking technology that transfers data directly between the memory of
two computers without involving the CPU or operating system of either host.
\textit{RoCE} and InfiniBand are common RDMA transports for storage.

\item[\textbf{RGW (RADOS Gateway)}]
The Ceph component that provides S3- and Swift-compatible object storage
APIs on top of \textit{RADOS}. RGW supports multi-site replication and
bucket-level policies.

\item[\textbf{RoCE (RDMA over Converged Ethernet)}]
A network protocol that allows \textit{RDMA} over standard Ethernet
infrastructure. RoCEv2 runs over UDP/IP and is widely used for
\textit{NVMe-oF} and GPU interconnects.

\item[\textbf{RPO (Recovery Point Objective)}]
The maximum acceptable age of data that can be lost after a failure,
measured from the last successful backup or replication. An RPO of zero
requires synchronous replication.

\item[\textbf{RTO (Recovery Time Objective)}]
The maximum acceptable duration from a failure event to full service
restoration. RTO drives the choice between hot standby, warm failover,
and cold restore strategies.

\end{description}

%% ---------------------------------------------------------------------------
\section*{S}
%% ---------------------------------------------------------------------------

\begin{description}

\item[\textbf{S3 (Simple Storage Service)}]
An object storage API originally defined by Amazon Web Services. S3 has
become a de facto standard: many on-premises systems (\textit{StorageGRID},
\textit{MinIO}, Ceph \textit{RGW}) implement S3-compatible endpoints.

\item[\textbf{SAN (Storage Area Network)}]
A dedicated high-speed network that provides block-level access to storage.
SANs typically use \textit{FC}, \textit{iSCSI}, or \textit{NVMe-oF}
transports and present \textit{LUNs} to connected hosts.

\item[\textbf{SCSI (Small Computer Systems Interface)}]
A set of standards for physically connecting and transferring data between
computers and peripheral devices. The SCSI command set underpins FC, iSCSI,
and SAS storage protocols.

\item[\textbf{SED (Self-Encrypting Drive)}]
A storage device with built-in hardware encryption that transparently
encrypts all data written to the media. SEDs conform to the \textit{TCG
Opal} standard and protect data at rest without performance penalty.

\item[\textbf{SELinux (Security-Enhanced Linux)}]
A mandatory access control (MAC) framework in the Linux kernel that enforces
security policies beyond traditional discretionary access controls.
SELinux labels can restrict which processes may access storage paths.

\item[\textbf{SLOG (Separate Log Device)}]
A dedicated device that provides persistent storage for the ZFS Intent Log
(\textit{ZIL}) in \textit{ZFS}. A fast SLOG (NVMe or Optane) accelerates
synchronous write workloads (e.g., NFS, databases).

\item[\textbf{SMB (Server Message Block)}]
A network file-sharing protocol used primarily in Windows environments. SMB
3.x adds encryption, multichannel, continuous availability, and RDMA
(\texttt{SMB Direct}) support. See also \textit{CIFS}.

\item[\textbf{SnapLock}]
An ONTAP compliance feature that provides tamper-proof \textit{WORM}
storage. SnapLock volumes prevent data modification or deletion until a
configurable retention period expires, meeting regulatory requirements such
as SEC 17a-4 and \textit{GDPR}.

\item[\textbf{SnapMirror}]
A NetApp replication technology that efficiently mirrors data between ONTAP
systems (on-premises or cloud) at the volume or consistency-group level.
SnapMirror supports asynchronous, synchronous, and semi-synchronous modes.

\item[\textbf{SnapVault}]
A NetApp disk-to-disk backup technology based on \textit{SnapMirror} that
maintains a vault of read-only snapshot copies for long-term retention.

\item[\textbf{StorageGRID}]
A NetApp software-defined object storage platform that provides an
S3-compatible interface with policy-driven data placement, erasure coding,
and information lifecycle management (ILM) across multiple sites.

\item[\textbf{SVM (Storage Virtual Machine)}]
A logical partition within an \textit{ONTAP} cluster that owns its own
volumes, \textit{LIFs}, and protocol configuration. SVMs enable secure
multi-tenancy on shared hardware.

\item[\textbf{SyncMirror}]
An ONTAP feature that synchronously mirrors data at the aggregate level
across two physical storage pools (plexes), providing local high
availability. SyncMirror is a building block of \textit{MetroCluster}.

\end{description}

%% ---------------------------------------------------------------------------
\section*{T}
%% ---------------------------------------------------------------------------

\begin{description}

\item[\textbf{TBW (Terabytes Written)}]
A cumulative SSD endurance metric representing the total amount of data
that can be written to a drive over its warranty life. See also
\textit{DWPD}.

\item[\textbf{TCG Opal}]
A specification from the Trusted Computing Group that defines a standard
interface for \textit{SED} management, including key management, locking
ranges, and authentication.

\item[\textbf{Trident}]
NetApp's open-source \textit{CSI} driver for Kubernetes. Trident provisions
persistent volumes on ONTAP (NFS, iSCSI, NVMe-oF), SolidFire, and cloud
storage back-ends. Now part of NetApp BlueXP.

\end{description}

%% ---------------------------------------------------------------------------
\section*{V}
%% ---------------------------------------------------------------------------

\begin{description}

\item[\textbf{VADP (vSphere APIs for Data Protection)}]
VMware APIs that enable backup applications to create and manage VM
snapshots, perform \textit{CBT}-based incremental backups, and restore
virtual machines without installing agents inside guests.

\item[\textbf{VDO (Virtual Data Optimizer)}]
A Linux kernel module (now integrated into LVM as \texttt{lvcreate --type
vdo}) that provides inline deduplication and compression at the block
layer.

\item[\textbf{VFS (Virtual File System)}]
An abstraction layer in the Linux kernel that provides a uniform interface
for user-space programs to interact with different filesystem
implementations (ext4, XFS, Btrfs, NFS, etc.).

\item[\textbf{vSAN}]
VMware's hyper-converged, software-defined storage layer that aggregates
local disks across ESXi hosts into a shared datastore with policy-driven
placement, erasure coding, and deduplication.

\end{description}

%% ---------------------------------------------------------------------------
\section*{W}
%% ---------------------------------------------------------------------------

\begin{description}

\item[\textbf{WAFL (Write Anywhere File Layout)}]
NetApp's proprietary copy-on-write filesystem at the heart of
\textit{ONTAP}. WAFL never overwrites data in place, enabling efficient
snapshots, clones, and consistency guarantees.

\item[\textbf{WekaFS}]
A high-performance, POSIX-compliant parallel filesystem designed for
GPU-intensive AI/ML workloads. WekaFS stripes data across NVMe tiers and
supports tiering to object storage.

\item[\textbf{WORM (Write Once, Read Many)}]
A storage policy that prevents modification or deletion of written data.
WORM is a regulatory requirement in financial services, healthcare, and
government. See also \textit{SnapLock}.

\item[\textbf{WWNN (World Wide Node Name)}]
A globally unique 64-bit identifier assigned to a Fibre Channel node (e.g.,
a storage controller or HBA). Each node may have multiple \textit{WWPN}
ports.

\item[\textbf{WWPN (World Wide Port Name)}]
A globally unique 64-bit identifier assigned to an individual port on a
Fibre Channel node. WWPNs are used in SAN zoning and LUN masking to
control host access.

\end{description}

%% ---------------------------------------------------------------------------
\section*{X}
%% ---------------------------------------------------------------------------

\begin{description}

\item[\textbf{XFS}]
A high-performance 64-bit journaling filesystem originally developed by SGI.
XFS excels at large files and parallel I/O, supports online growth, and is
the default filesystem on RHEL, CentOS, and Rocky Linux.

\end{description}

%% ---------------------------------------------------------------------------
\section*{Z}
%% ---------------------------------------------------------------------------

\begin{description}

\item[\textbf{ZFS (Zettabyte File System)}]
A combined volume manager and filesystem originally developed by Sun
Microsystems. ZFS provides copy-on-write, end-to-end checksumming,
snapshots, clones, inline compression, deduplication, \textit{RAID-Z}, and
\textit{dRAID}. On Linux it is available through OpenZFS.

\item[\textbf{ZIL (ZFS Intent Log)}]
A write-ahead log in \textit{ZFS} that records synchronous write
transactions. The ZIL is replayed after a crash to guarantee transactional
consistency. A dedicated \textit{SLOG} device can accelerate ZIL commits.

\item[\textbf{ZNS (Zoned Namespace SSDs)}]
An NVMe specification where the drive's address space is divided into
sequential-write-required zones. ZNS reduces write amplification, improves
endurance, and lowers per-GB cost by eliminating the device-side Flash
Translation Layer for those zones.

\end{description}
